\section{Discussion}
\label{sec:discussion}

\subsection{What is established}

The answer to the uniform-gap question at the top of the pendulum is negative: no constant $c > 0$ independent of $\sigma$ gives $\kappa - R \geq c\sigma$ for all states. The counterexamples are numerical, but every number was confirmed by evaluations that share no code, and their behaviour, $g \to 1/(2k_0)$ with $k_0$ limited to $\sqrt{\ln(1/\sigma)}$ by the cost of the blur, is explained by the same mechanism that the proofs use.

What survives is a logarithmic law, and both of its sides have the same form. The strips show that $g_{\min}(\sigma)$ can decrease like $1/\sqrt{\ln(1/\sigma)}$. The lower bounds show that it cannot decrease faster: for every state with a bounded ratio $K_\Lambda$ (Theorem~\ref{thm:log}), and without any ratio condition for product states in any orientation (Theorem~\ref{thm:product}) and for states without deep valleys (Proposition~\ref{prop:novalley}), the class that contains the strips. The states that come closest to the edge are therefore covered by a theorem without any ratio condition, whose one hypothesis, the absence of deep valleys, we verified on them numerically. On the strips the best proved constant is below the numerical one by a factor of about 7 at the smallest $\sigma$, and most of the factor comes from two identified steps. This is the familiar structure of a limit theorem with a matching construction, here with the limit proved for broad classes and the construction established numerically.

The proofs also isolate what controls the gap: the single number $\xb$, the fineness of the structure along the unstable direction after the blur. In observer language, the gap is a tug of war between the information along the stable direction and the nonlinearity felt where the state is sharp (Theorem~\ref{thm:direct}). Fine structure lets a state be sharp close to the top, which is what a small gap requires, but the blur erases fine structure exponentially (Lemmas~\ref{lem:smoothing1d} and~\ref{lem:split}), and a state confined near the top cannot afford to lose its last structure along the stable direction. The result is a ceiling $\xb \lesssim 2\ln(1/\sigma)$ and a gap of order $\sigma/\sqrt{\ln(1/\sigma)}$. The second route adds a second picture: following the blur as it grows, the only thing that can make the observer lose information faster is a valley, a place where the observer cannot tell from which part of the state a point came, and such a valley sits in a dip whose slopes carry information (Theorem~\ref{thm:valley}).

The effect matters only at very fine resolution. Gaussian mixtures with up to 32 components stayed above $0.884\,\sigma$, the best Gaussian state tends to $\sigma$~\cite{namba2026gap}, and the strips beat the mixtures only below $\sigma \approx 7\times 10^{-6}$. At such resolutions the factor $1/\sqrt{\ln(1/\sigma)}$ is close to $0.3$. For practical purposes the resolution gap of Refs.~\cite{namba2026gap,namba2026many} stands as a proportionality to $\sigma$; as a mathematical statement it carries a logarithm.

\subsection{Open problems}

\begin{enumerate}
\item \emph{The logarithmic law for all states.} No counterexample to Conjecture~\ref{conj:general} is known. The intermediate inequality~\eqref{eq:T2} holds on all 1778 stored pairs of a state and a blur, with $tR \leq \ln(1/(tJ_{ww})) + 6.54$, but not with universal constants for every state: hidden sources of alternating sign in the valley between very long lines defeat it (Sec.~\ref{sec:removing}). A proof for general states must therefore use the confinement (H2), which caps $\ln(1/(\sigma^2J_{ww}))$ at $2\ln(1/\sigma) + O(1)$ through Eq.~\eqref{eq:CR}, and not only the structure of the blur. We shall report on this separately.
\item \emph{The remainders.} $\varepsilon$, $\varepsilon_m$ and $\varepsilon_\pi$ are explicit and were small on every focused state ($|\varepsilon| \leq 0.004$, $\varepsilon_m \leq 0.11$), but state-independent bounds on $\varepsilon$ are not known.
\item \emph{The strips analytically.} The limit $g \to 1/(2k_0)$, the law~\eqref{eq:fit} and the absence of deep valleys are numerical. An asymptotic analysis of Eq.~\eqref{eq:strip} would turn the construction side into a theorem.
\item \emph{Beyond the pendulum.} The mechanism uses only the Gaussian blur, the uncertainty principle and the confinement of the state near a hyperbolic point. We expect a logarithmic law at the hyperbolic points of other flows and in many degrees of freedom, with the matrix formula of Ref.~\cite{namba2026many} in place of the pendulum's constant, but this has not been checked.
\end{enumerate}

\subsection{Relation to other work}

Finite-resolution entropies are used as diagnostics of chaos, where the growth rate of a coarse-grained entropy tracks the Lyapunov exponent~\cite{kannan2026}. Those studies fix the resolution; here the resolution widens in time, and the question is the margin by which the onset of entropy growth stays below the stretching rate. The heat-flow part of Sec.~\ref{sec:removing} belongs to the study of Fisher information along the heat flow, where de~Bruijn's identity and the $\Gamma_2$ calculus are standard tools~\cite{cover2006,bakry2014} and convexity properties are known for log-concave densities~\cite{toscani2015}; the directional and multimodal setting here, and the role of valleys, appear to be outside those results. For phase-space diffusion models~\cite{dockhorn2022,namba2026threshold,namba2026learned}, the present result says that data with fine structure can narrow the band between the two thresholds $R$ and $\kappa$ found in Ref.~\cite{namba2026learned}, and, for the classes covered by the theorems, only logarithmically in the noise level.
